\begin{ex}
Tính: $35\text{ kg} + 28\text{ kg}$. Kết quả là bao nhiêu ki-lô-gam?
\shortans{63}
\loigiai{
Ta có: $35 + 28 = 63$ kg.
}
\end{ex}

\begin{ex}
Tính: $72\text{ kg} - 35\text{ kg}$. Kết quả là bao nhiêu ki-lô-gam?
\shortans{37}
\loigiai{
Ta có: $72 - 35 = 37$ kg.
}
\end{ex}

\begin{ex}
Tính: $46\text{ }l + 27\text{ }l$. Kết quả là bao nhiêu lít?
\shortans{73}
\loigiai{
Ta có: $46 + 27 = 73$ lít.
}
\end{ex}

\begin{ex}
Tính: $81\text{ }l - 45\text{ }l$. Kết quả là bao nhiêu lít?
\shortans{36}
\loigiai{
Ta có: $81 - 45 = 36$ lít.
}
\end{ex}

\begin{ex}
Một ngày có bao nhiêu giờ?
\shortans{24}
\loigiai{
Một ngày có 24 giờ.
}
\end{ex}

\begin{ex}
Một giờ có bao nhiêu phút?
\shortans{60}
\loigiai{
Một giờ có 60 phút.
}
\end{ex}

\begin{ex}
Tháng 5 có bao nhiêu ngày?
\shortans{31}
\loigiai{
Tháng 5 có 31 ngày.
}
\end{ex}

\begin{ex}
Tháng 6 có bao nhiêu ngày?
\shortans{30}
\loigiai{
Tháng 6 có 30 ngày.
}
\end{ex}

\begin{ex}
Một tuần lễ có bao nhiêu ngày?
\shortans{7}
\loigiai{
Có 7 ngày.
}
\end{ex}

\begin{ex}
Bao gạo nặng 45 kg, bao ngô nặng hơn bao gạo 8 kg. Bao ngô nặng bao nhiêu ki-lô-gam?
\shortans{53}
\loigiai{
Bao ngô nặng: $45 + 8 = 53$ kg.
}
\end{ex}

\begin{ex}
Thùng dầu có 65 lít, người ta lấy ra 28 lít. Thùng còn lại bao nhiêu lít dầu?
\shortans{37}
\loigiai{
Còn lại: $65 - 28 = 37$ lít.
}
\end{ex}

\begin{ex}
Đồng hồ chỉ 9 giờ 30 phút. Kim dài chỉ vào số mấy?
\shortans{6}
\loigiai{
Kim dài chỉ số 6.
}
\end{ex}

\begin{ex}
Đồng hồ chỉ 14 giờ. Đây là mấy giờ chiều?
\shortans{2}
\loigiai{
Là 2 giờ chiều.
}
\end{ex}

\begin{ex}
Đoạn dây thứ nhất dài 38 cm, đoạn thứ hai dài 25 cm. Cả hai đoạn dài bao nhiêu xăng-ti-mét?
\shortans{63}
\loigiai{
Dài: $38 + 25 = 63$ cm.
}
\end{ex}

\begin{ex}
Sợi dây đồng dài 90 cm, cắt bớt 35 cm. Sợi dây còn lại dài bao nhiêu xăng-ti-mét?
\shortans{55}
\loigiai{
Còn lại: $90 - 35 = 55$ cm.
}
\end{ex}

\begin{ex}
Bé Nam nặng 26 kg, khi bé bế thêm con mèo nặng 3 kg thì cả bé và mèo nặng bao nhiêu ki-lô-gam?
\shortans{29}
\loigiai{
Nặng: $26 + 3 = 29$ kg.
}
\end{ex}

\begin{ex}
Can thứ nhất có 15 lít nước, can thứ hai có 18 lít nước. Cả hai can có bao nhiêu lít nước?
\shortans{33}
\loigiai{
Có: $15 + 18 = 33$ lít.
}
\end{ex}

\begin{ex}
Nếu Thứ Sáu là ngày 14 thì Thứ Sáu tuần sau là ngày bao nhiêu?
\shortans{21}
\loigiai{
Ngày đó là: $14 + 7 = 21$.
}
\end{ex}

\begin{ex}
Nếu Thứ Ba là ngày 24 thì Thứ Ba tuần trước là ngày bao nhiêu?
\shortans{17}
\loigiai{
Ngày đó là: $24 - 7 = 17$.
}
\end{ex}

\begin{ex}
Tính giá trị biểu thức: $24\text{ kg} + 18\text{ kg} - 15\text{ kg}$.
\shortans{27}
\loigiai{
Ta có $24 + 18 = 42$, rồi $42 - 15 = 27$ kg.
}
\end{ex}

\begin{ex}
Tính giá trị biểu thức: $50\text{ }l - 25\text{ }l + 18\text{ }l$.
\shortans{43}
\loigiai{
Ta có $50 - 25 = 25$, rồi $25 + 18 = 43$ lít.
}
\end{ex}

\begin{ex}
Một lớp học có buổi dã ngoại kéo dài từ 8 giờ sáng đến 3 giờ chiều (15 giờ). Buổi dã ngoại kéo dài bao nhiêu giờ?
\shortans{7}
\loigiai{
Thời gian: $15 - 8 = 7$ giờ.
}
\end{ex}

\begin{ex}
Đĩa cân bên trái có gói đường và quả cân 1 kg, đĩa bên phải có quả cân 5 kg thì cân thăng bằng. Gói đường nặng bao nhiêu ki-lô-gam?
\shortans{4}
\loigiai{
Gói đường nặng: $5 - 1 = 4$ kg.
}
\end{ex}

\begin{ex}
Có 5 can nước, mỗi can chứa 3 lít. Cả 5 can chứa bao nhiêu lít nước?
\shortans{15}
\loigiai{
Chứa: $3 \times 5 = 15$ lít.
}
\end{ex}

\begin{ex}
Một đoạn đường dài 100 mét, công nhân đã sửa được 65 mét. Còn lại bao nhiêu mét đường chưa sửa?
\shortans{35}
\loigiai{
Còn lại: $100 - 65 = 35$ mét.
}
\end{ex}
